% !TEX program = pdflatex
\documentclass[11pt]{article}
\usepackage{series}
\usepackage{booktabs}
\usepackage{tabularx}
\usepackage{array}
\usepackage{xcolor}

\newcommand{\Z}{\mathbb{Z}}
\newcommand{\C}{\mathbb{C}}
\newcommand{\Pic}{\operatorname{Pic}}
\newcommand{\Inv}{\operatorname{Inv}}
\newcommand{\Hom}{\operatorname{Hom}}
\newcommand{\Det}{\operatorname{Det}}
\newcommand{\Pf}{\operatorname{Pf}}
\newcommand{\cOne}{c_1}
\newcommand{\one}{\mathbf{1}}
\newcommand{\Line}{\mathcal{L}}
\newcommand{\Coeff}{\mathcal{C}}
\newcommand{\Fields}{\mathcal{F}}
\newcommand{\Ops}{\mathcal{O}}
\newcommand{\Bkg}{\mathcal{B}}
\newcommand{\GaugeGrp}{\mathcal{G}}
\newcolumntype{Y}{>{\raggedright\arraybackslash}X}
\newcolumntype{P}[1]{>{\raggedright\arraybackslash}p{#1}}
\hypersetup{
  pdftitle={Bundle Selection Rules, Anomaly Lines, and Charge Lattices},
  pdfauthor={Peter Nero}
}

\title{\textbf{Bundle Selection Rules, Anomaly Lines, and Charge Lattices}\\
\large A Conditional Topological Layer for Modal Triplet Theory}
\author{Peter Nero}
\date{Version 2, July 2026}

\begin{document}
\maketitle

\begin{abstract}
Topology can constrain a four-dimensional effective theory, but only after
the global gauge group, field representations, internal bundles, zero-mode
spaces, coefficient bundles, and contraction maps have been specified. This
paper replaces an earlier topology-only argument by that typed statement.
For a fixed global gauge group, characters of its \(U(1)\) factor form an
integer lattice; rational hypercharges arise after a normalization choice and
quotient compatibility conditions. This lattice does not by itself select
the observed matter representations. Gauge anomalies are encoded by the
determinant or Pfaffian line of a family of chiral Dirac operators over
background-field space, rather than by the determinant of the matter bundle
on spacetime. Local anomaly is detected by curvature and global anomaly by
holonomy; cancellation requires a compatible local equivariant
trivialization. For effective operators, we give the exact tensor-product
line classes of the Yukawa, Weinberg, \(QQQL\),
\(u^cu^cd^ce^c\), and singlet-Majorana monomials. The two baryon-number
violating examples are Standard Model gauge singlets and, in nonsupersymmetric
SMEFT, have mass dimension six. They are excluded only if a declared
realization supplies an additional bundle, symmetry, cohomology, or overlap
obstruction. We prove a sufficient bundle-selection theorem and explain why
passing its tests does not guarantee a nonzero coupling. Current Modal
Triplet Theory (MTT) has exact finite results for a selected chiral
representation, its anomaly table, the faithful
\((SU(3)\times SU(2)\times U(1))/\Z_6\) group, and a unique anomaly-free
shared hypercharge direction within the chosen finite completion. Those
results do not yet select the physical compactification endpoint or prove
that every dangerous operator is absent. The result is a rigorous
realization-by-realization selection framework, not a universal topology-only
derivation of the Standard Model.
\end{abstract}

\section*{Version 2 revision note}

\begin{description}[style=nextline]
\item[Supersedes.]
Version 1, DOI \href{https://doi.org/10.5281/zenodo.18261774}
{\texttt{10.5281/zenodo.18261774}}.
\item[Reason.]
Version 1 used rational tensor powers of a line bundle without first choosing
a root, placed the anomaly line on the wrong base, treated existence of a
global section as equivalent to bundle triviality, and suggested that
Standard Model gauge topology removes operators that are already Standard
Model gauge singlets.
\item[Resolution.]
Version 2 separates character lattices, global gauge-group descent,
determinant-line anomalies, internal Picard classes, coefficient spaces, and
actual overlap pairings. Every displayed operator now has an explicit class
and a stated decision rule.
\item[Retained content.]
The useful core survives: topology can provide exact, pre-dynamical
obstructions inside a fixed realization, and those obstructions are valuable
early consistency and falsifiability tests.
\item[Open boundary.]
MTT has not yet selected the physical bundle endpoint and coefficient
functional that would decide all proton-decay and lepton-number-violating
operators. Its current exact finite gauge and anomaly results are recorded
without promoting them to a no-knob compactification theorem.
\end{description}

\tableofcontents

\section{What a topology-only claim must contain}

The phrase ``forbidden by topology'' sounds stronger than it usually is. It
can mean at least four different things:
\begin{enumerate}
\item a proposed field is not a representation of the global gauge group;
\item a family of chiral fermions has an uncancelled local or global anomaly;
\item an internal bundle product has no allowed coefficient or invariant
pairing;
\item a particular overlap integral vanishes.
\end{enumerate}
These statements live on different spaces and require different proofs. A
charge is a representation label. An anomaly is a property of a quantum
fermion theory over its background-field space. An internal selection rule
is a statement about bundles and mode spaces on a compactification or
overlap base. A vanishing coefficient is a statement about a specific
multilinear functional. Conflating them makes an argument look shorter, but
also makes its conclusion undecidable.

\subsection{The typed realization record}

Let \(X\) denote the internal or overlap space relevant to a proposed
four-dimensional realization. It need not be physical space and it is not
identified with a spatial slice. Let
\[
G=\frac{SU(3)\times SU(2)\times U(1)}{\Gamma}
\]
be the declared global gauge group, where \(\Gamma\) is a specified central
subgroup. For every four-dimensional field label \(f\), a usable record
contains:
\[
\bigl(R_f,n_f,\Line_f,V_f,\nabla_f\bigr).
\]
Here \(R_f\) is the nonabelian representation, \(n_f\in\Z\) is the normalized
\(U(1)\) character, \(\Line_f\to X\) is the internal line or more general
bundle factor, \(V_f\) is the selected zero-mode space, and \(\nabla_f\) is
the connection when holonomy or parallel transport matters.

For a monomial
\[
\Ops=\prod_{a=1}^{r} f_a,
\]
define
\begin{align}
R_{\Ops}&=\bigotimes_{a=1}^{r}R_{f_a},&
n_{\Ops}&=\sum_{a=1}^{r}n_{f_a},&
\Line_{\Ops}&=\bigotimes_{a=1}^{r}\Line_{f_a}.
\label{eq:operator-data}
\end{align}
Dual or conjugate fields contribute the dual representation, opposite
character, and dual line bundle.

\begin{definition}[Operator completion record]\label{def:completion}
An operator completion record consists of the data in
\eqref{eq:operator-data}, an allowed coefficient space
\[
\Coeff_{\Ops}\subseteq
\Gamma\!\left(X,\Line_{\Ops}^{*}\right),
\]
a \(G\)-invariant contraction of the nonabelian representations, and a
multilinear functional
\[
\mu_{\Ops}:
\Coeff_{\Ops}\otimes V_{f_1}\otimes\cdots\otimes V_{f_r}
\longrightarrow \C.
\]
The four-dimensional coefficient is the value of \(\mu_{\Ops}\) on the
selected coefficient and zero modes.
\end{definition}

This definition accommodates several familiar cases. A constant scalar
coefficient corresponds to a chosen trivialization of
\(\Line_{\Ops}\). A holomorphic compactification may instead use a section of
\(\Line_{\Ops}^{*}\) and a cohomological cup product. A Kaluza--Klein
reduction may use an integral of mode representatives. A discrete symmetry
can set \(\Coeff_{\Ops}=0\). The final scalar cannot be inferred from the
symbol \(\Line_{\Ops}\) alone.

\subsection{Four distinctions that prevent false selection rules}

\begin{proposition}[Section, trivialization, and connection]\label{prop:distinctions}
Let \(\Line\to X\) be a complex line bundle.
\begin{enumerate}
\item A nowhere-zero global section trivializes \(\Line\), but an arbitrary
global section may have zeros and need not trivialize it.
\item If \(\cOne(\Line)\ne0\), then \(\Line\) is not topologically trivial.
The converse can fail when torsion or flat Picard data are present.
\item A flat connection can have nontrivial holonomy.
\item A trivial line bundle can carry a connection with nonzero curvature.
\end{enumerate}
A nonzero global parallel section exists exactly when the connection has
trivial holonomy; such a section supplies a connection-preserving
trivialization.
\end{proposition}

\begin{proof}
The first statement follows by using a nowhere-zero section as a global
frame. The Chern class obstructs a topological trivialization, but does not
classify every line bundle on every space. Flatness sets the curvature to
zero but leaves a representation of \(\pi_1(X)\) as holonomy. Conversely, on
the trivial line, \(d+iA\) has curvature \(i\,dA\), which need not vanish.
Parallel transport of a nonzero vector produces a path-independent global
parallel section exactly when all holonomies fix it; for a line this means
trivial holonomy.
\end{proof}

\begin{remark}
Version 1 repeatedly moved between these notions. In the revised argument,
``topologically trivial,'' ``flat,'' ``trivial holonomy,'' and ``equipped
with a selected trivialization'' are never synonyms.
\end{remark}

\section{Charge lattices and what they do not select}

\subsection{Integer characters come first}

The continuous characters of \(U(1)\) are
\[
\chi_n(z)=z^n,\qquad n\in\Z.
\]
Thus a principal \(U(1)\) bundle \(P_Y\) and its weight-one associated line
\(K\) generate the charged lines
\[
K_n=P_Y\times_{\chi_n}\C\cong K^{\otimes n},
\qquad
\cOne(K_n)=n\,\cOne(K).
\]
There is no canonical rational tensor power \(K^{\otimes q}\). Writing a
rational physical hypercharge \(Y=n/N\) presupposes a normalization \(N\),
or equivalently a primitive line whose integer weight is \(n\). If one starts
from a nonprimitive line and wants an \(N\)-th root, existence of that root
is additional global data.

\begin{theorem}[Fixed-group charge lattice]\label{thm:charge-lattice}
Fix a compact global gauge group
\[
G=(G_{\mathrm{ss}}\times U(1))/\Gamma
\]
with finite central \(\Gamma\). Every finite-dimensional representation of
\(G\) pulls back to a representation of \(G_{\mathrm{ss}}\times U(1)\) with
integer \(U(1)\) character \(n\). It descends to \(G\) exactly when the
combined action of every element of \(\Gamma\) is trivial. Consequently the
allowed labels form an integer lattice with congruence conditions. After a
normalization \(Y=n/N\), they form a rational lattice. This determines
allowed representation labels, not which labels occur in nature.
\end{theorem}

\begin{proof}
Every continuous character of \(U(1)\) is \(\chi_n\) for a unique
\(n\in\Z\). A representation of the quotient is precisely a representation
of the covering product on which the quotient subgroup \(\Gamma\) acts
trivially. The descent condition is therefore a finite set of congruences
relating \(n\) to the central characters of \(G_{\mathrm{ss}}\). Selecting a
matter spectrum is extra data: the representation ring contains many
allowed elements.
\end{proof}

For the conventional left-handed Standard Model labels, the useful integer
normalization is \(n=6Y\):
\begin{table}[ht]
\centering
\caption{Conventional one-family labels. These are an example to be
reproduced, not a consequence of the abstract lattice theorem.}
\label{tab:sm-labels}
\begin{tabular}{@{}ccccccc@{}}
\toprule
field & \(Q\) & \(u^c\) & \(d^c\) & \(L\) & \(e^c\) & \(N^c\) \\
\midrule
\(SU(3)\times SU(2)\) &
\((\mathbf3,\mathbf2)\) &
\((\bar{\mathbf3},\mathbf1)\) &
\((\bar{\mathbf3},\mathbf1)\) &
\((\mathbf1,\mathbf2)\) &
\((\mathbf1,\mathbf1)\) &
\((\mathbf1,\mathbf1)\) \\
\(6Y\) & \(1\) & \(-4\) & \(2\) & \(-3\) & \(6\) & \(0\) \\
\bottomrule
\end{tabular}
\end{table}

The global form of the gauge group matters. Different quotients of groups
with the same Lie algebra can impose different descent conditions
\cite{Hashimoto2013}. This is why ``integral cohomology quantizes
hypercharge'' is incomplete: one must name the primitive character, the
global group, its quotient, and the field representations.

\section{Anomalies live over background-field space}

\subsection{The determinant line has a different base}

Let \(\Bkg\) be a space of gauge and metric backgrounds and let
\(\GaugeGrp\) be the relevant gauge group. A chiral fermion representation
defines a family of chiral Dirac operators
\[
D_b^+:\Gamma(S^+\otimes E_b)\longrightarrow
\Gamma(S^-\otimes E_b),
\qquad b\in\Bkg.
\]
Its determinant line is the family-index line
\[
\Det D^+\longrightarrow \Bkg
\quad\text{or, equivariantly, over }\Bkg/\GaugeGrp.
\]
It is not the ordinary top exterior power \(\det E_b\) of the matter bundle
over spacetime. The Bismut--Freed connection on this family line has a
curvature determined by the local index density, while its holonomy detects
global anomaly information \cite{Freed1995,Freed2014}.

\begin{theorem}[Scoped anomaly criterion]\label{thm:anomaly}
For a declared family of chiral Dirac operators, a nonzero equivariant
curvature of the determinant or Pfaffian line is a local anomaly
obstruction, and nontrivial equivariant holonomy of a flat anomaly line is a
global anomaly obstruction. Anomaly cancellation on that domain requires a
local, gauge-compatible trivialization of the anomaly theory; in the
determinant-line model this includes vanishing curvature and holonomy after
all allowed counterterms and inflow contributions are included.
\end{theorem}

\begin{proof}
The fermionic functional integral is naturally a section of the determinant
or Pfaffian line. A gauge-compatible scalar partition function requires a
compatible trivialization. Curvature obstructs local flatness and encodes
the perturbative anomaly. If curvature vanishes, nontrivial holonomy still
prevents a single-valued gauge-compatible section around loops in
background-field space. Conversely, a compatible local equivariant
trivialization converts the section into a scalar functional on the stated
domain. The locality qualifier is essential: an arbitrary nonlocal
trivialization is not an admissible counterterm.
\end{proof}

\subsection{The conventional one-family arithmetic}

The representation labels in \cref{tab:sm-labels} satisfy the familiar
four-dimensional anomaly equations. Omitting common positive normalization
factors,
\begin{align}
\mathcal A_{SU(3)^2Y}
 &=2Y_Q+Y_{u^c}+Y_{d^c}=0,\\
\mathcal A_{SU(2)^2Y}
 &=3Y_Q+Y_L=0,\\
\mathcal A_{\mathrm{grav}^2Y}
 &=6Y_Q+3Y_{u^c}+3Y_{d^c}+2Y_L+Y_{e^c}=0,\\
\mathcal A_{Y^3}
 &=6Y_Q^3+3Y_{u^c}^3+3Y_{d^c}^3+2Y_L^3+Y_{e^c}^3=0.
\end{align}
The color cubic anomaly cancels between two triplet components of \(Q\) and
the two antitriplets \(u^c,d^c\). The number of weak doublets per family,
counting color, is \(3+1=4\), so the Witten \(SU(2)\) global anomaly also
cancels. These equations verify a chosen representation. They do not by
themselves derive that representation or its multiplicity.

\section{Exact operator classes and selection rules}

\subsection{Field lines}

Attach internal line bundles
\[
\Line_Q,\ \Line_u,\ \Line_d,\ \Line_L,\ \Line_e,\
\Line_N,\ \Line_H
\quad\text{in }\Pic(X)
\]
to the left-handed fields in \cref{tab:sm-labels} and to a Higgs doublet
\(H\) with \(6Y_H=3\). Write
\[
\ell_f=[\Line_f]\in\Pic(X)
\]
additively. Complex conjugation contributes \(-\ell_f\). The following table
now makes every claim checkable.

\begin{table}[p]
\centering
\small
\caption{Gauge-neutral monomials and their internal line classes. A zero
class is necessary for a constant coefficient but is not, by itself,
sufficient for a nonzero coupling.}
\label{tab:operators}
\begin{tabularx}{\textwidth}{@{}lcccY@{}}
\toprule
monomial & 4D dimension & \(6Y\) sum & internal Picard class
& exact additional test \\
\midrule
\(QHu^c\) &
\(4\) &
\(1+3-4=0\) &
\(\ell_Q+\ell_H+\ell_u\) &
An \(SU(3)\times SU(2)\) singlet exists. A constant Yukawa needs a selected
trivialization; a geometric Yukawa needs a nonzero allowed overlap
functional.\\
\addlinespace
\(QH^\dagger d^c\) &
\(4\) &
\(1-3+2=0\) &
\(\ell_Q-\ell_H+\ell_d\) &
Same distinction between a trivial line and a nonzero overlap coefficient.\\
\addlinespace
\(LH^\dagger e^c\) &
\(4\) &
\(-3-3+6=0\) &
\(\ell_L-\ell_H+\ell_e\) &
Gauge neutrality does not determine the charged-lepton Yukawa matrix.\\
\addlinespace
\((LH)(LH)\) &
\(5\) &
\(2(-3+3)=0\) &
\(2\ell_L+2\ell_H\) &
The Weinberg coefficient lies in the dual line and a symmetric flavor
pairing; it is not decided by a bare ``real structure'' statement.\\
\addlinespace
\(QQQL\) &
\(6\) &
\(3(1)-3=0\) &
\(3\ell_Q+\ell_L\) &
It is already an SM gauge singlet. Exclusion requires this extra class,
an additional symmetry, a vanishing cohomology group, or a zero overlap
functional.\\
\addlinespace
\(u^cu^cd^ce^c\) &
\(6\) &
\(2(-4)+2+6=0\) &
\(2\ell_u+\ell_d+\ell_e\) &
It is also an SM gauge singlet. The same realization-specific certificate
is required.\\
\addlinespace
\(N^cN^c\) &
\(3\) &
\(0\) &
\(2\ell_N\) &
A bare Majorana mass needs an allowed invariant symmetric bilinear or a
coefficient section of \(\Line_N^{-2}\).\\
\bottomrule
\end{tabularx}
\end{table}

In nonsupersymmetric SMEFT, \(QQQL\) and \(u^cu^cd^ce^c\) are four-fermion
operators of mass dimension six \cite{Grzadkowski2010}. In supersymmetric
language, analogous four-chiral-superfield superpotential monomials are
often called dimension-five proton-decay operators after the superspace
measure and effective suppression are accounted for. The framework must
state which convention it uses.

\subsection{The selection theorem}

\begin{theorem}[Typed bundle-selection rule]\label{thm:selection}
Fix an operator completion record as in \cref{def:completion}. The operator
\(\Ops\) is absent on the declared realization if any one of the following
holds:
\begin{enumerate}
\item \(\Inv_G(R_{\Ops})=0\);
\item the allowed coefficient space \(\Coeff_{\Ops}\) is zero;
\item every allowed invariant multilinear functional
\(\mu_{\Ops}\) vanishes on the selected zero-mode spaces.
\end{enumerate}
If coefficients are restricted to constants, nontriviality of
\(\Line_{\Ops}\) is a sufficient obstruction. In particular,
\(\cOne(\Line_{\Ops})\ne0\) is a sufficient certificate in that restricted
constant-coefficient setting. Passing all of these tests is necessary for a
nonzero coupling but does not guarantee one.
\end{theorem}

\begin{proof}
Without a \(G\)-invariant contraction, the monomial cannot be a gauge
scalar. Without an allowed dual coefficient, its internal bundle factor
cannot be paired into a scalar. If every allowed multilinear functional
vanishes, dimensional reduction emits zero coefficient. For a constant
coefficient, pairing requires a selected trivialization of
\(\Line_{\Ops}\); a nontrivial line has none. A nonzero first Chern class
certifies nontriviality. Conversely, a gauge singlet with a trivial total
line can still have zero coefficient because its zero-mode product,
symmetry selection rule, or overlap integral vanishes.
\end{proof}

\begin{corollary}[What must be shown for the two proton-decay examples]
\label{cor:proton}
Topology removes \(QQQL\) or \(u^cu^cd^ce^c\) from a chosen realization only
after that realization proves, respectively,
\[
3\ell_Q+\ell_L\ne0,
\qquad
2\ell_u+\ell_d+\ell_e\ne0,
\]
in the relevant Picard or differential-cohomology group under a
constant-coefficient rule, or supplies an equivalent vanishing certificate
for the allowed coefficient space or overlap functional.
\end{corollary}

\begin{remark}
This is weaker than Version 1 rhetorically and stronger mathematically. It
does not announce proton stability from generic overlap language. It tells a
calculation exactly what class or functional must be emitted before that
claim can be made. Concrete heterotic line-bundle models illustrate this
model-by-model logic: additional \(U(1)\) symmetries can forbid particular
operators, but the allowed spectrum is explicitly computed for each bundle
choice \cite{Anderson2012}.
\end{remark}

\section{Majorana terms: the precise condition}

A ``real structure'' is one sufficient route to a Majorana pairing, but it
is not the most precise universal criterion. Let \(E\) be the gauge and
internal bundle carried by a left-handed Weyl field. A bare mass requires a
gauge-invariant symmetric bilinear in the internal/flavor labels compatible
with the antisymmetric spinor contraction. Equivalently, the relevant
singlet must occur in the correct symmetry component of \(E\otimes E\), and
its coefficient bundle must be allowed. A real representation often
provides such a bilinear; a pseudoreal representation may require multiple
flavors. A complex charged representation generally does not allow a bare
Majorana term without symmetry breaking or an additional field.

For a sterile \(N^c\), the internal line condition is the last row of
\cref{tab:operators}. For active Standard Model neutrinos, the leading
gauge-invariant route is the dimension-five Weinberg operator
\((LH)(LH)\) \cite{Weinberg1979}. The exact internal class is
\[
2\ell_L+2\ell_H.
\]
If the operator is absent at one level, higher-dimensional operators,
symmetry-breaking insertions, or a different completion can still generate
a Majorana mass. Therefore failure of a bare real structure does not imply
that every neutrino mass must be Dirac.

\section{What current MTT has actually established}

The current corpus is stronger than the abstract Version 1 picture in some
finite directions and weaker in the physical compactification direction.
The distinction is important enough to state row by row.

\begin{table}[ht]
\centering
\small
\caption{Current MTT status relevant to this paper. ``Exact'' refers to the
declared finite carrier, not automatic selection of the physical vacuum.}
\label{tab:mtt-status}
\begin{tabularx}{\textwidth}{@{}P{0.19\textwidth}P{0.18\textwidth}Y@{}}
\toprule
object & current tier & statement \\
\midrule
Typed family carrier (A46) &
derived exact; vacuum selector open &
A 48-state family-diagonal chiral representation emits the \(Q,u^c,d^c,
L,e^c,N^c\) rows and exactly verifies the local and global anomaly table.\\
\addlinespace
Native gauge group (A47) &
derived exact &
The selected rank-\(1,2,3\) finite carriers have native
\(U(1),SU(2),SU(3)\) automorphisms, and the diagonal center kernel is
\(\Z_6\). The faithful group is
\((SU(3)\times SU(2)\times U(1))/\Z_6\).\\
\addlinespace
Shared hypercharge line (A50) &
derived exact; profile tier &
Within the chosen finite completion, the abelian anomaly equations have the
unique primitive null vector that yields
\(6Y=(1,-4,2,-3,6,0)\). This adds no continuous knob, but it is a
low-energy finite-profile theorem.\\
\addlinespace
\(E_6\) \(Q_\psi\) color anomaly (A22) &
derived exact; strong-CP map open &
Three matter families contribute \(+12\), complete-\(\mathbf{27}\) exotics
contribute \(-12\), and the selected trace cancels. The axion-current
threshold map remains a separate obligation.\\
\addlinespace
Embedded renormalized-SM equivalence (A04) &
closed at declared profile standard &
The twelve-obligation baseline is closed at one shared physical primitive
and measured-profile standard. This is not zero-knob prediction.\\
\addlinespace
Strict no-knob upgrade (A05) &
open, \(2/9\) closed &
Physical source selection, remaining values, and stronger prediction
claims are not supplied by the topology-only layer.\\
\bottomrule
\end{tabularx}
\end{table}

The exact finite results are reproducible in the curated MTT calculation
repository \cite{MTTResults} and summarized in the current typed Standard
Model compatibility and noncommutative-geometry papers
\cite{MTTB4,MTTNCG}. They improve the charge and anomaly part of the old
paper substantially. They do not yet provide the operator-specific internal
classes in \cref{tab:operators} on a selected physical visible--hidden
compactification. In particular, no current theorem may replace the two
conditions in \cref{cor:proton} by the word ``generic.''

\section{Falsifiability after the correction}

The revised framework produces sharp tests, but each test is conditional on
a declared realization.

\begin{enumerate}[label=\textbf{F\arabic*.}]
\item \textbf{Global-group descent.}
If a claimed field representation fails the \(\Gamma\)-descent condition,
the realization is inconsistent.
\item \textbf{Anomaly line.}
If the complete fermion and inflow content leaves nonzero anomaly curvature
or holonomy, the quantum gauge theory is inconsistent on that domain.
\item \textbf{Operator certificate.}
If a paper claims that an operator is absent, it must publish
\(R_{\Ops},n_{\Ops},\Line_{\Ops},\Coeff_{\Ops},V_f,\mu_{\Ops}\) or an
equivalent complete certificate. A nonzero observed coefficient contradicts
that declared certificate.
\item \textbf{Charge selection.}
An allowed charge lattice is not a prediction of the observed assignment.
A predictive claim must also identify the selection theorem that chooses
the observed primitive vector and matter multiplicities.
\item \textbf{Held-out consequence.}
Selection data used to construct a model cannot count again as a prediction.
A physical claim requires an observable not used to choose the bundles,
representations, or coefficient functional.
\end{enumerate}

An observed proton-decay channel would therefore falsify an MTT realization
that had already certified the corresponding operator as absent. It would
not, by itself, falsify the abstract idea that some other overlap
realizations possess topological selection rules. This distinction turns a
broad slogan into a testable scientific statement.

\section{Version 2 changes and reasons}

\begin{table}[ht]
\centering
\small
\begin{tabularx}{\textwidth}{@{}P{0.28\textwidth}YY@{}}
\toprule
Version 1 statement & Version 2 decision & Reason \\
\midrule
Rational charges use \(L_Y^{\otimes q}\). &
Replace by integer characters \(K^{\otimes n}\), quotient descent, and an
explicit normalization \(Y=n/N\). &
Rational tensor powers are not defined without root data.\\
\addlinespace
Integral cohomology selects hypercharge. &
Narrow to an allowed charge lattice; record the separate A50 finite
selection theorem. &
A lattice does not select a spectrum or normalization by itself.\\
\addlinespace
The anomaly line is \(\det\mathcal E\to\Sigma\). &
Replace by \(\Det D^+\to\Bkg/\GaugeGrp\) with curvature, holonomy, locality,
counterterm, and inflow data. &
The fermionic anomaly is a family-index object over background-field space.\\
\addlinespace
An operator is allowed iff its bundle has a global section. &
Replace by the completion record and \cref{thm:selection}. &
A section may vanish; nontrivial bundles may have sections; a coupling
requires an invariant scalar functional.\\
\addlinespace
\(QQQL\) and \(u^cu^cd^ce^c\) are leading dimension-five operators
forbidden by overlap topology. &
Correct to nonsupersymmetric dimension six, show both are SM gauge singlets,
and state their exact extra Picard classes. &
Their exclusion is realization-specific, not a consequence of SM
hypercharge.\\
\addlinespace
Majorana mass requires a real structure. &
Use the invariant symmetric bilinear and coefficient-space criterion; keep
real structure as one sufficient route. &
Representation type, flavor multiplicity, Higgs insertions, and higher
operators all matter.\\
\addlinespace
Topology-only results explain proton stability universally. &
Withdraw. &
Current MTT has not emitted the two physical operator-class certificates.\\
\bottomrule
\end{tabularx}
\end{table}

\section{Conclusion}

Topology is most useful here as a compiler of exact obstructions. Once a
global gauge group and a physical bundle realization are fixed, it can say
that a representation does not descend, an anomaly line cannot be
trivialized, or an operator has no allowed coefficient or overlap
functional. Those are strong, pre-dynamical conclusions.

Topology cannot, without additional source data, choose the observed matter
spectrum, make every dangerous Standard Model gauge singlet disappear, or
turn a vanishing first Chern class into a nonzero coupling. Version 2 makes
that boundary explicit. The present MTT finite carrier already supplies an
exact \(\Z_6\) global group, a chiral anomaly table, and a unique
anomaly-free normalized hypercharge vector inside its chosen completion.
The next physical theorem is now concrete: emit the selected internal
classes and coefficient functionals for the operator rows in
\cref{tab:operators} on the same compactification branch. Until then, the
paper provides a rigorous conditional selection language and an auditable
completion contract.

\bibliographystyle{plain}
\bibliography{main}

\end{document}
